\section{A positive repair for Puiseux data: a separated simple branch}
\label{sec:simplebranch}

The Puiseux packet in the repository identifies a precise obstruction:
arbitrary Cauchy presentations of the coefficients do not uniformly supply
individually accessible exact coefficients of a Puiseux root. Its example
\(Y^2-X^2-a\) makes the distinction between \(a=0\) and \(a>0\) visible in
one returned coefficient. Finite exact algebraic presentations avoid that
particular problem, but they are not the only usable input interface.
Here we establish a second interface: a supplied root of the constant
polynomial together with a finite certificate separating its derivative
from zero. The coefficients may otherwise be arbitrary primitive
recursive complex names. The result gives one ordinary power-series
branch, with ramification integer \(1\), explicit coefficient evaluators,
and an explicit disk of convergence. It does not address the singular
case where that derivative vanishes.

The proof has a deliberately finite core. Every estimate first concerns
finite polynomials with Gaussian-rational coefficients. Cauchy names are
introduced only after those estimates have been proved. This order makes
the quantifier-free induction requirements visible and avoids using a
fixed-point theorem or analytic completeness as an arithmetic axiom.
The repository motivation is
\texttt{algebra/puiseux\_ipra/evidence/scope\_obstruction.md} in the
proof collection \cite{ProveIt}.

\subsection{The finite polynomial kernel}

For a Gaussian rational \(z=u+iv\), put
\(\lVert z\rVert_1=|u|+|v|\). This norm is submultiplicative:
expanding the two coordinates of a product and applying the rational
triangle inequality gives
\(\lVert zw\rVert_1\leq\lVert z\rVert_1\lVert w\rVert_1\).
For a finite polynomial \(U(X)=\sum_k u_kX^k\) and a positive rational
\(r\), define its weighted coefficient norm by
\[
  \lVert U\rVert_r=\sum_k\lVert u_k\rVert_1r^k.
\]
Finite convolution gives
\[
 \lVert U+V\rVert_r\leq\lVert U\rVert_r+\lVert V\rVert_r,
 \qquad
 \lVert UV\rVert_r\leq\lVert U\rVert_r\lVert V\rVert_r.
 \tag{A1}\label{eq:branch-norm}
\]
These are inequalities between rational numbers. To prove the product
inequality, expand the coefficient at each degree, use
\(r^{a+b}=r^ar^b\), and rearrange the resulting finite double sum.
Induction on the lengths of the arrays proves each rearrangement.
There is no assertion about an infinite sum in (A1).

Let
\[
 Q(X,Z)=\sum_{i+j\leq D}q_{ij}X^iZ^j,
 \qquad q_{00}=q_{01}=0,
 \qquad D\geq2,
 \qquad \sum_{i,j}\lVert q_{ij}\rVert_1\leq B,
\]
where \(B\geq1\) is rational. Define
\[
 s=\frac{1}{4DB},\qquad
 r=\frac{s}{4B}=\frac{1}{16DB^2}.
 \tag{A2}\label{eq:branch-radii}
\]
In particular, \(0<r\leq s\leq1\). All zero coefficients in this
description are specified positions in a finite array. No procedure is
asked to decide whether a represented coefficient vanishes.

\begin{lemma}[A finite contraction estimate]
\label{lem:finite-poly-contraction}
If \(\lVert U\rVert_r,\lVert V\rVert_r\leq s\), then
\[
 \lVert Q(X,U)-Q(X,V)\rVert_r
       \leq\tfrac14\lVert U-V\rVert_r,
 \qquad
 \lVert Q(X,U)\rVert_r\leq s/2.
 \tag{A3}\label{eq:finite-poly-contraction}
\]
\end{lemma}
\begin{proof}
For every positive integer \(j\), the polynomial identity
\[
 U^j-V^j=(U-V)\sum_{h<j}U^hV^{j-1-h}
\]
follows by cancellation of consecutive terms. Its weighted norm is at
most \(j s^{j-1}\lVert U-V\rVert_r\). In a term of \(Q\) that depends
on \(Z\), the excluded position \((0,1)\) ensures
\(i+j-1\geq1\). Consequently
\[
 r^i s^{j-1}\leq s,
 \quad
 \sum_{i\geq0,\,j\geq1}j\lVert q_{ij}\rVert_1r^i s^{j-1}
 \leq DBs=\tfrac14.
\]
Here the sum is over terms with \(j\geq1\); the index \(i\) may be
zero. This proves the first inequality. The other excluded position,
\((0,0)\), gives
\(\lVert Q(X,0)\rVert_r\leq Br=s/4\). Applying the first inequality
with \(V=0\) gives
\(\lVert Q(X,U)\rVert_r\leq s/4+s/4=s/2\).
\end{proof}

Define the finite polynomial iterates
\[
 U_0=0,\qquad U_{n+1}=Q(X,U_n).
 \tag{A4}\label{eq:finite-iterates}
\]
Their degrees can be bounded by \(D^n\), so they can be constructed by
primitive recursion on finite coefficient arrays. If only coefficients
through degree \(K\) are required, truncating after every operation to
that degree gives exactly the same retained coefficients. Truncation
discards only terms of higher degree, and every exponent in \(Q\) is
nonnegative.

\begin{lemma}[Stabilization by degree]
\label{lem:branch-stabilization}
Every \(U_n\) has zero constant coefficient and
\(\lVert U_n\rVert_r\leq s\). Moreover
\[
  X^{n+1}\mid U_{n+1}-U_n.
 \tag{A5}\label{eq:degree-stabilization}
\]
Thus, if \(c_k=[X^k]U_k\) for \(k\geq1\), the coefficient of
\(X^k\) in every \(U_N\), \(N\geq k\), is \(c_k\).
\end{lemma}
\begin{proof}
The zero constant coefficient and the norm bound follow by induction,
using \(q_{00}=0\), \(U_0=0\), and (A3). Suppose \(U,V\) have zero
constant coefficient and \(X^n\mid U-V\). The factorization used in
the preceding proof shows that
\(Q(X,U)-Q(X,V)=(U-V)R(X)\). Every term of \(R\) has a factor
\(X\): if \(i\geq1\), that factor is explicit; if \(i=0\), then
\(j\geq2\), and each product in
\(\sum_{h<j}U^hV^{j-1-h}\) contains at least one factor having zero
constant coefficient. Hence \(X^{n+1}\mid Q(X,U)-Q(X,V)\).
The initial difference \(U_1-U_0=Q(X,0)\) is divisible by \(X\).
Induction gives (A5). Summing the finitely many later differences
proves the assertion about each retained coefficient.
\end{proof}

Put \(S_K(X)=\sum_{1\leq k\leq K}c_kX^k\), with \(S_0=0\).

\begin{theorem}[Finite majorant, tail, and residual]
\label{thm:finite-branch}
For every \(K\leq N\),
\[
 \sum_{1\leq k\leq N}\lVert c_k\rVert_1r^k\leq s.
 \tag{A6}\label{eq:branch-majorant}
\]
If \(\lVert x\rVert_1\leq r/2\), then
\[
 \begin{split}
 \lVert S_N(x)-S_K(x)\rVert_1&\leq s\,2^{-(K+1)},\\
 \lVert Q(x,S_K(x))-S_K(x)\rVert_1&\leq s\,2^{-K}.
 \end{split}
 \tag{A7}\label{eq:branch-tail}
\]
The polynomial \(Q(X,S_K)-S_K\) has zero coefficients through
degree \(K\).
\end{theorem}
\begin{proof}
By stabilization, \(S_N\) is the truncation of \(U_N\), so (A6)
follows from its norm bound. For the first estimate in (A7), expand the
finite tail. For \(k>K\),
\(\lVert x\rVert_1^k\leq r^k2^{-k}\leq r^k2^{-(K+1)}\).
Summing and using (A6) gives the claimed bound.

The coefficients through degree \(K\) of \(Q(X,S_K)\) agree with
those of \(Q(X,U_K)=U_{K+1}\). Stabilization identifies these with
the coefficients of \(S_K\). Write
\(R_K=Q(X,S_K)-S_K\). Its low coefficients therefore vanish, while
(A3) gives
\(\lVert R_K\rVert_r\leq s/2+s\leq2s\). Every remaining degree
is at least \(K+1\). The same finite-tail estimate gives
\[
 \lVert R_K(x)\rVert_1
 \leq2^{-(K+1)}\lVert R_K\rVert_r
 \leq s2^{-K}.
\]
This proves all claims without introducing an infinite series.
\end{proof}

\subsection{The precision loss is explicit and independent of iteration count}

For the application to names, a perturbation estimate is essential.
Let \(\widetilde Q\) have the same excluded coefficient positions,
degree bound, and majorant \(B\), and suppose
\(\lVert q_{ij}-\widetilde q_{ij}\rVert_1\leq\delta\) at every
position. Set \(m=(D+1)^2\), a convenient upper bound on the number
of positions. For \(\lVert U\rVert_r\leq s\),
\[
 \lVert Q(X,U)-\widetilde Q(X,U)\rVert_r\leq m\delta,
 \tag{A8}\label{eq:coefficient-perturbation}
\]
because \(r,s\leq1\). Applying (A3) to the two sequences of
iterates yields, with \(E_n=\lVert U_n-\widetilde U_n\rVert_r\),
\[
 E_0=0,\qquad E_{n+1}\leq E_n/4+m\delta,
 \qquad E_n\leq2m\delta.
 \tag{A9}\label{eq:iterate-perturbation}
\]
The last inequality is proved by induction: its right-hand side is
preserved because \(2m\delta/4+m\delta\leq2m\delta\).
Thus the dependence on the coefficient input does not get worse as the
number of iterations grows.

Let \(e_p=2^{-p}\). A regular complex name in this subsection is a
primitive recursive Gaussian-rational function \(a(p)\) satisfying
\[
 \lVert a(p)-a(q)\rVert_1\leq e_p+e_q.
 \tag{A10}\label{eq:branch-name}
\]
We work with a fixed primitive recursive family of such functions; this
is a schema, not a universal evaluator for all primitive recursive
programs. Assume each coefficient \(q_{ij}\) is supplied in this form,
with the two excluded entries identically zero. A rational majorant
valid at \emph{every} precision is available directly from the input:
one may take
\[
 B=\max\left(1,\sum_{i+j\leq D}
                 (\lVert q_{ij}(0)\rVert_1+2)\right).
 \tag{A11}\label{eq:uniform-majorant}
\]
Indeed (A10) bounds \(\lVert q_{ij}(p)\rVert_1\) by
\(\lVert q_{ij}(0)\rVert_1+1+e_p\leq\lVert q_{ij}(0)\rVert_1+2\).
A sharper certified majorant can be used instead. This distinction
matters: the finite proof now applies literally to the rational
polynomial \(Q_p\) at each precision, without an informal passage from
exact coefficient inequalities to floating-point ones.

For a nonnegative rational \(u\), write \(\lambda(u)\) for the least
\(h\leq\lceil u\rceil+1\) such that
\(2^h\geq\max(1,u)\). This is a bounded primitive recursive search;
\(2^h\geq h+1\) proves that the indicated bound suffices. Let
\(U_k^p\) be the \(k\)-th iterate for \(Q_p\). An explicit name for
the \(k\)-th coefficient is
\[
 \begin{split}
 h_k&=\lambda(2m r^{-k}),\\
 p(k,n)&=n+h_k+2,\\
 C(k,n)&=[X^k]U_k^{p(k,n)}.
 \end{split}
 \tag{A12}\label{eq:coefficient-name}
\]
For \(k=0\), set \(C(0,n)=0\). The coefficient estimate from
(A9) gives
\[
 \lVert [X^k]U_k^p-[X^k]U_k^q\rVert_1
 \leq2m r^{-k}(e_p+e_q).
\]
Substituting (A12) bounds the difference between \(C(k,n)\) and
\(C(k,n')\) by \((e_n+e_{n'})/4\). Therefore these are regular
names, uniformly in the numerical index \(k\). All loops may be
truncated at degree \(k\).

There is also a uniform evaluator for the finite partial sums. Suppose
\(\lVert x\rVert_1,\lVert x'\rVert_1\leq r/2\), and
\(\lVert U\rVert_r\leq s\). Factoring \(x^k-(x')^k\) gives
\[
 \lVert U(x)-U(x')\rVert_1
 \leq \frac{s}{r}\lVert x-x'\rVert_1.
 \tag{A13}\label{eq:polynomial-eval-lipschitz}
\]
To check the constant, the \(k\)-th summand is bounded by
\(\lVert u_k\rVert_1 k(r/2)^{k-1}\lVert x-x'\rVert_1\),
and \(k2^{-(k-1)}\leq1\) for every \(k\geq1\).

Let a name \(x(p)\) satisfy both (A10) and
\(\lVert x(p)\rVert_1\leq r/2\) at every precision. Put
\[
 T_K(p)=\bigl(\text{truncation of }U_K^p
                    \text{ through degree }K\bigr)(x(p)),
 \qquad C_* =2m+s/r.
\]
Combining (A9), truncation, and (A13) gives the entirely finite estimate
\[
 \lVert T_K(p)-T_K(q)\rVert_1
       \leq C_*(e_p+e_q),
 \tag{A14}\label{eq:partial-sum-name}
\]
with no dependence of \(C_*\) on \(K\).
The bounded-name presentation of the disk is harmless. For a rational
radius \(R>0\), define the rational radial projection by
\(\pi_R(z)=z\) when \(\lVert z\rVert_1\leq R\), and
\(\pi_R(z)=Rz/\lVert z\rVert_1\) otherwise. It has Lipschitz
constant at most \(2\). To see this, assume
\(\lVert z\rVert_1\geq\lVert w\rVert_1\). If both are outside,
subtract the two scaled vectors, keep \(R/\lVert z\rVert_1\) as
the common factor, and bound the remaining scale difference using
\(\lVert z\rVert_1-\lVert w\rVert_1\leq\lVert z-w\rVert_1\).
If only \(z\) is outside, the same calculation uses
\(\lVert z\rVert_1-R\leq\lVert z\rVert_1-\lVert w\rVert_1\).
If both are inside, the claim is immediate. Thus applying \(\pi_R\)
to a regular name at index \(n+1\) gives a regular name in the disk.
It represents the original value whenever that value is in the disk,
because \(\pi_R\) fixes that value and the same Lipschitz estimate
controls approximation errors. Alternatively, one can require the
bounded presentation directly in the input contract.

If \(y_0\) is a further regular name, define
\[
 \begin{split}
 K(n)&=n+3,\\
 p(n)&=n+\lambda(C_*+1)+3,\\
 Y(n)&=y_0(p(n))+T_{K(n)}(p(n)).
 \end{split}
 \tag{A15}\label{eq:root-name}
\]
This is a regular name. For example, if \(n\leq n'\), compare
\(T_{K(n)}(p(n))\) first with \(T_{K(n)}(p(n'))\), using (A14),
and then with \(T_{K(n')}(p(n'))\), using the finite tail in (A7)
for the one rational polynomial \(Q_{p(n')}\). The sum of these
errors and the error in \(y_0\) is at most
\[
 (C_*+1)(e_{p(n)}+e_{p(n')})
          +s2^{-(n+4)}
 \leq\tfrac18(e_n+e_{n'})+\tfrac1{16}e_n
 \leq e_n+e_{n'}.
\]
The case \(n'\leq n\) is symmetric and is decided by natural-number
comparison. This proof constructs the name; it does not extract it from
an assertion that an unspecified sequence converges.

\subsection{Normalization and inversion from one finite separation certificate}

We now explain how the excluded entries of \(Q\) arise without an exact
zero test. Let
\(P(X,Y)=\sum_{i+j\leq D}a_{ij}X^iY^j\), let \(y_0\) be a
supplied regular name, and assume the supplied arithmetic contracts
prove \(P(0,y_0)=0\). Let
\(d=P_Y(0,y_0)\). The additional input is a finite derivative
separation witness. One convenient form consists of \(t\) and a
coordinate \(j\in\{\operatorname{Re},\operatorname{Im}\}\) such that
\[
       |d_j(t)|\geq4e_t.
 \tag{A16}\label{eq:finite-apartness}
\]
This is one decidable rational test, supplied as data; it is not a
conclusion drawn from the negation of \(d=0\).

For completeness, here is an explicit inverse construction. For every
\(p\geq t+2\), regularity gives
\[
 |d_j(p)|\geq4e_t-e_t-e_p\geq\tfrac{11}{4}e_t>2e_t.
\]
Choose a rational \(M\geq1\) bounding
\(\lVert d(t)\rVert_1+2\). Then
\(\lVert d(p)\rVert_1\leq M\) at every precision, and, at the
specified precisions,
\[
 |d(p)|_2^2\geq4e_t^2,
 \qquad
 \lVert d(p)^{-1}\rVert_1\leq
 H:=\frac{M}{4e_t^2}.
\]
The rational inverse is
\((u-iv)/(u^2+v^2)\). Its denominator is strictly positive by the
displayed bound. The identity
\(z^{-1}-w^{-1}=z^{-1}(w-z)w^{-1}\) gives
\[
 \lVert d(p)^{-1}-d(q)^{-1}\rVert_1
             \leq H^2(e_p+e_q).
\]
Consequently
\[
 I(n)=d\bigl(\max(t+2,n+\lambda(H^2)+2)\bigr)^{-1}
 \tag{A17}\label{eq:inverse-name}
\]
is a regular name. The exact rational identity
\(d(p)d(p)^{-1}=1\), together with regularity and the magnitude
bound, proves that it represents an inverse of \(d\). This proof uses
only finite rational identities and explicit error estimates.

Expand the following finite expression:
\[
 Q(X,Z)=Z-d^{-1}P(X,y_0+Z).
 \tag{A18}\label{eq:normalized-branch}
\]
The constant coefficient is zero by the supplied root equation, and
the coefficient at \(Z\) is \(1-d^{-1}d=0\). In the resulting
coefficient array, put these two positions identically equal to the
rational zero. The finite binomial expansion and the supplied equalities
prove that this modified array represents (A18). All other coefficient
names are obtained by finite sums, products, and the inverse (A17).
Here are explicit schedules for those operations. If names \(a,b\)
have rational bounds \(A,B\) at all precisions, addition uses
\(a(n+1)+b(n+1)\), with bound \(A+B\). Multiplication uses
\(a(p)b(p)\), where \(p=n+\lambda(A+B)+1\), with bound \(AB\).
Indeed the product difference is at most
\((A+B)(e_p+e_q)\). A finite expression is evaluated recursively
using these shifts at every node; (A17) handles its only division.
This gives a primitive recursive evaluator and a magnitude bound for
each normalized coefficient. This construction
uses a proof of the two zero identities; it never branches on a test of
equality of arbitrary complex names.

\begin{theorem}[One simple analytic branch with a primitive recursive name]
\label{thm:simple-branch}
For every fixed primitive recursive family of the preceding input data,
the finite certificate (A16), the supplied equality \(P(0,y_0)=0\), and
the regular-name contracts yield the following primitive recursive
outputs and arithmetic contracts.
There is a family of coefficient names \(c_k\), explicitly given by
(A12), with \(c_0=0\), and the rational radii (A2), such that:
\begin{enumerate}
\item For every finite degree \(N\), substitution of
      \(y_0+\sum_{k=1}^{N}c_kX^k\) into \(P\) has zero coefficients
      through degree \(N\).
\item The majorant (A6) holds for every finite coefficient prefix.
      On the disk \(\lVert x\rVert_1\leq r/2\), the finite tail is
      bounded by \(s2^{-(K+1)}\).
\item If \(M\) is a rational bound on \(\lVert d\rVert_1\), then
      \[
       \left\lVert P\left(x,y_0+S_K(x)\right)\right\rVert_1
          \leq Ms2^{-K}.
       \tag{A19}\label{eq:original-residual}
      \]
\item The evaluator (A15) gives a regular root name \(Y(x)\).
      It satisfies \(P(x,Y(x))=0\) in the represented equality
      relation, and \(\lVert Y(x)-y_0\rVert_1\leq s\).
\item It is the unique represented root in that ball.
\end{enumerate}
Thus the output is an ordinary power-series branch; its ramification
integer is \(e=1\).
\end{theorem}
\begin{proof}
The finite polynomial results apply to each rational approximation
\(Q_p\), since the uniform majorant (A11) was chosen for every
precision. The perturbation estimates identify the coefficient names
defined in (A12) and prove that every finite coefficient prefix satisfies
the asserted identities and bounds. In elementary real or complex
notation, this says that \(S_K\) has the stabilized coefficients
established in Lemma~\ref{lem:branch-stabilization}. In arithmetic
notation, each assertion says that its rational approximations differ
by at most the stated error at every requested precision.

The finite low-coefficient identity for \(Q(X,S_N)-S_N\), multiplied
by \(-d\), proves the first claim, using (A18). Multiplication of
the second estimate in (A7) by the magnitude bound \(M\) proves
(A19). Formula (A15) and its finite error calculation give the root
name. To make the last error budget explicit, let \(A_{ij}\) bound
the norms of the original coefficient names at every precision and
let \(A_0\) bound the name of \(y_0\). Put
\[
 R_x=1,\qquad R_y=A_0+s+1,\qquad
 L=\sum_{i+j\leq D,\,j\geq1}
       jA_{ij}R_x^iR_y^{j-1}.
\]
Factoring differences of powers proves
\(\lVert P(x,u)-P(x,v)\rVert_1\leq L\lVert u-v\rVert_1\)
on these bounded disks. The tail and (A19) therefore give
\[
 \lVert P(x,Y(x))\rVert_1\leq s(M+L)2^{-K}
\]
at every finite cutoff \(K\), in the represented inequality
notation. For a requested error \(e_n\), take
\[
 K=n+\lambda\bigl(s(M+L)+1\bigr)+2.
 \tag{A19a}\label{eq:final-zero-budget}
\]
The right-hand side is at most \(e_n/4\). Evaluate the finite
polynomial expressions and their equality certificates with error at
most \(e_n/4\), using the recursive addition, multiplication, and
inverse schedules just specified. This proves the rational zero
certificate with room to spare. All choices are primitive recursive.
The notation of a represented inequality is eliminated by exactly this
requested-precision argument, rather than by an axiom about limits.

Finally, suppose \(z,w\) are two roots after subtracting \(y_0\),
with norms at most \(s\). Use the radial projection described above
to choose equivalent names whose approximants lie in that ball.
The identity (A18) makes them fixed points
of \(Q(x,\cdot)\). The scalar version of (A3) gives
\(\lVert z-w\rVert_1\leq\frac14\lVert z-w\rVert_1\), hence
\(\lVert z-w\rVert_1=0\). At finite precision the same calculation
has an additional known error \(\eta\) and gives
\(\lVert z-w\rVert_1\leq4\eta/3\). Choosing the common input
precision so that this is below \(e_n\) proves the represented
equality. No test of whether either norm is zero occurs.
\end{proof}

\paragraph{What is internal to iPRA.}
The theorem is a schema for fixed evaluator families, with ordinary
numerical parameters for the degree, coefficient indices, precision,
and separation witness. Gaussian rationals, arrays, binomial
coefficients, convolutions, truncations, and finite sums have primitive
recursive codes. The core of (A3)--(A9) is already a collection of
rational inequalities. For a requested final cutoff, all array indices
are bounded by explicitly constructed sizes. Induction on a running
iteration count proves a quantifier-free invariant containing the norm
bound and the bounded list of zero low coefficients. The perturbation
invariant is \(E_n\leq2m\delta\), a rational inequality. Its precision
indices are free numerical parameters during induction. Consequently
no induction formula quantifies over all precisions, all complex
numbers, or all functions. Universally quantified contracts are invoked
only at their finitely many needed numerical instances.

The notation involving represented values in the theorem is an
explanation of these finite contracts. Formalizing those contracts as
\(\mathrm{Prf}\) objects in the repository remains work to be done; a
Lean theorem about complete complex numbers would not establish that
object-language conclusion.

\subsection{Two worked branches and the remaining singular boundary}

For \(P(X,Y)=Y^2-1-X\) and \(y_0=1\), one has \(d=2\) and
\[
 Q(X,Z)=(X-Z^2)/2.
\]
Take \(D=2\), \(B=1\), \(s=1/8\), \(r=1/32\). The explicitly
certified evaluation disk is \(\lVert x\rVert_1\leq1/64\). The
first stabilized coefficients are
\[
 c_1=\frac12,\quad c_2=-\frac18,\quad c_3=\frac1{16},
 \quad c_4=-\frac5{128},\quad c_5=\frac7{256}.
\]
They are obtained by finite polynomial iteration, without selecting a
complex square root at each step. The root is the unique one in
\(\lVert Y-1\rVert_1\leq1/8\), and
\[
 \lVert P(x,1+S_K(x))\rVert_1\leq\tfrac14\,2^{-K}.
\]
The radius is intentionally conservative; the theorem prioritizes a
uniform finite proof over a sharp analytic radius.

For \(P(X,Y)=Y+Y^3-X\), take \(y_0=0\), \(d=1\),
\(Q(X,Z)=X-Z^3\), \(D=3\), and \(B=2\). The construction gives
\(s=1/24\), \(r=1/192\), and coefficients beginning
\[
  c_1=1,\quad c_2=0,\quad c_3=-1,\quad c_4=0,
  \quad c_5=3,\quad c_6=0,\quad c_7=-12.
\]
This is a concrete local inversion theorem with a fully specified
primitive recursive precision schedule.

Return now to the obstruction polynomial \(Y^2-X^2-a\). If a root
\(y_0\) of \(Y^2-a\) is supplied together with a separation witness
for \(2y_0\), the theorem gives its simple branch and exact
coefficient names. When \(a=0\), that witness does not exist for
\(y_0=0\), although the two branches \(Y=\pm X\) can be written
down in that special case. There is no uniform rule here for deciding
which case occurs. The theorem therefore resolves the regular branch
construction for Cauchy-presented coefficients while preserving the
repository's genuine obstruction at the singular transition.

\subsection{Finite mesh certificates replace one common use of compactness}
\label{sec:mesh-certificates}

The preceding result uses a supplied quantitative hypothesis. A related
repair applies whenever an argument invokes a positive minimum on a
compact domain. A finite mesh and a continuity modulus can supply the
needed lower bound, provided the finite numerical tests have enough
margin. They do not manufacture that margin from the bare assertion
that a function has no zeros.

\begin{proposition}[A finite certificate of positivity]
\label{prop:mesh-positivity}
Let \(f\) be a fixed real evaluator on a rational box
\([0,1]^d\), with an approximation procedure \(A(q,p)\) satisfying
\(|A(q,p)-f(q)|\leq e_p\) in the represented sense. Suppose a
primitive recursive modulus \(\mu\) certifies
\[
 \lVert x-y\rVert_\infty\leq2^{-\mu(k)}
       \quad\Longrightarrow\quad
 |f(x)-f(y)|\leq2^{-k}.
\]
Fix a positive rational \(\gamma\). Choose \(k,p\) so that
\(2^{-k}\leq\gamma\) and \(e_p\leq\gamma\), and let
\(M=2^{\mu(k)}\). If the finite tests
\[
 A(j_1/M,\ldots,j_d/M,p)\geq3\gamma
       \quad(0\leq j_i\leq M)
 \tag{A20}\label{eq:mesh-certificate}
\]
all hold, then \(f(x)\geq\gamma\) throughout the box.
\end{proposition}
\begin{proof}
For a rational point of the box, choose in each coordinate a nearest
grid point, or simply the lower grid point, with the endpoint \(1\)
handled by the last grid index. The sup-norm distance is at most
\(1/M\). Each tested grid value of \(f\) is at least
\(3\gamma-e_p\geq2\gamma\). The modulus then gives
\(f(x)\geq2\gamma-2^{-k}\geq\gamma\).
The grid indices and the selected point are obtained by bounded
integer division. The test (A20) is a bounded conjunction, and the
argument is a rational inequality after denominator clearing.
For a represented point, first clamp each coordinate of its rational
approximants to \([0,1]\); this operation is rational and does not
increase their distance from any point of the box. One can then apply
this rational theorem to the resulting approximants and the supplied modulus,
or search the finite grid against one approximant with an additional
precision margin. There is no need to decide which exact grid cell
contains an arbitrary real point.
\end{proof}

For a complex function, apply the proposition to a certified real
quantity such as \(|g|_2^2\). Polynomial data supply the continuity
modulus by finite difference estimates. A successful mesh certificate
then supplies the denominator bound required for inversion, root
separation, or a contour argument. An unsuccessful certificate proves
nothing about the presence of a zero. In particular, increasing the
mesh until a test succeeds is not justified as a primitive recursive
algorithm without a separately supplied success bound.

The same method yields a rational upper certificate for a supremum:
the largest grid approximation plus its evaluation error plus the
modulus error bounds every function value. A grid point attaining the
finite maximum is an approximate maximizer. This supplies the
quantitative consequence often used from an extreme-value theorem,
while keeping the required continuity and precision data explicit.
